\begin{circuitikz}[
        american,
        voltage dir=RP,
        >={Latex[length=5mm, width=2mm]},
        %>=latex,
        alt={Schematic diagram of a simple DC electrical circuit connected to a PLC input point labeled Input Point. On the left, a DC battery labeled V provides voltage to the circuit. The upper conductor running directly from the positive terminal to the PLC module is labeled Common. The PLC input point is shown on the right inside a dashed red box, labeled NPN at the bottom, containing an internal current source component. The return path from the PLC module contains a push button switch labeled PB before connecting back to the negative terminal of the battery.}
    ]

    \draw[
        BakerBlack
    ]
    (0,0)
    to[battery, color=BakerBlack, v={$V$}] ++(0,2)
    to[short, color=BakerBlack, -o] coordinate[pos=0.5](com) ++(3,0)
    %coordinate[pos=0.5](com)
    coordinate(PLC_top)
    -- ++(1,0)
    to[isource, color=BakerBlack] ++(0,-2)
    to[short, color=BakerBlack, -o] ++(-1,0)
    coordinate(PLC_bottom)
    to[push button, color=BakerBlack, mirror, l_={PB}, font=\small] (0,0);

    \node[
        text=BakerBlack,
        font=\small,
        anchor=south east
    ]
    at (com)
    {Common};

    \node[
        draw=BakerADARed,
        dashed,
        anchor=north west,
        minimum width=2cm,
        minimum height=3cm,
        label={above:{\small \textcolor{BakerBlack}{Input Point}}},
    ]
    (PLC)
    at ($(PLC_top)+(0,0.5)$)
    {};

    \node[
        text=BakerBlack,
        anchor=south
    ]
    at (PLC.south)
    {{\tiny NPN}};

\end{circuitikz}